% Created 2023-03-02 jeu. 11:48
% Intended LaTeX compiler: pdflatex
\documentclass[presentation]{beamer}
\usepackage[utf8]{inputenc}
\usepackage[T1]{fontenc}
\usepackage{graphicx}
\usepackage{longtable}
\usepackage{wrapfig}
\usepackage{rotating}
\usepackage[normalem]{ulem}
\usepackage{amsmath}
\usepackage{amssymb}
\usepackage{capt-of}
\usepackage{hyperref}
\usepackage[french]{babel}
\usetheme{Frankfurt}
\usecolortheme{}
\usefonttheme{}
\useinnertheme{}
\useoutertheme{}
\author{Maxime Gaide, David Marcheix, Agnès Arnould, Xavier Skapin, Hakim Belhaouari, Stéphane Jean}
\date{\today}
\title{Détection statique des modifications topologiques dans les opérations de modélisation géométrique}
\setbeamertemplate{navigation symbols}{}
\hypersetup{
 pdfauthor={Maxime Gaide, David Marcheix, Agnès Arnould, Xavier Skapin, Hakim Belhaouari, Stéphane Jean},
 pdftitle={Détection statique des modifications topologiques dans les opérations de modélisation géométrique},
 pdfkeywords={},
 pdfsubject={},
 pdfcreator={Emacs 28.2 (Org mode 9.6)}, 
 pdflang={French}}
\begin{document}

\maketitle
\begin{frame}{Outline}
\tableofcontents
\end{frame}


\begin{frame}[label={sec:org9b85531}]{La détection de changement topologiques}
a
\end{frame}
\begin{frame}[label={sec:orgf8ac39a}]{Problématique de la détection dynamique}
b
\end{frame}
\begin{frame}[label={sec:orgb8d56e3}]{Méthode statique de détection}
c
\end{frame}
\begin{frame}[label={sec:org4c47204}]{Comparaison des méthodes statiques et dynamiques}
d
\end{frame}
\end{document}